% Coverage: physical pp. 95--97, Section 16.3 from its heading through the final sentence before Section 16.4; Eqs. (16.3.1)--(16.3.15) and the unnumbered convexity display.
% Corrected erratum: the constant segment has nonnegative, not nonpositive, second derivative, consistently with the convexity argument above.
\subsection{Energy Interpretation}

The effective action $\Gamma[\phi]$ and potential $V[\phi]$ have an important interpretation in terms of the energy and energy density respectively~\hyperref[ch16-ref-4]{[4]}. To see this, suppose that we turn on a current $\mathcal J^n(\mathbf{x},t)$ that rises smoothly from zero at $t=-\infty$ to a finite value $\mathcal J^n(\mathbf{x})$, and remains at that value for a long time $T$, after which it drops smoothly again to zero at $t=+\infty$. The effect of this perturbation is to convert the vacuum to a state with a definite energy $E[\mathcal J]$ (a functional of $\mathcal J^n(\mathbf{x})$), in which it remains for a time $T$, after which it returns to the vacuum. However, although the `out' vacuum is the same physical state as the `in' vacuum, the state vectors differ by the phase $\exp(-iE[\mathcal J]T)$ accumulated during the time $T$:
\begin{equation}
{}_{\mathrm{out}}\!\braket{\Omega}{\Omega}_{\mathrm{in},\mathcal J}
=\exp\bigl(-iE[\mathcal J]T\bigr).
\tag{16.3.1}\label{eq:16.3.1}
\end{equation}
Comparing with Eqs.~\eqref{eq:16.1.1} and \eqref{eq:16.1.3}, this gives
\begin{equation}
W[J]=-E[\mathcal J]T.
\tag{16.3.2}\label{eq:16.3.2}
\end{equation}

To see the connection between this energy and the effective action, suppose we seek the state $\ket{\Omega_\phi}$ that minimizes the energy expectation value
\begin{equation}
\langle H\rangle_\Omega
=\frac{\langle\Omega|H|\Omega\rangle}{\langle\Omega|\Omega\rangle},
\tag{16.3.3}\label{eq:16.3.3}
\end{equation}
subject to the condition that the quantum fields $\Phi_n(\mathbf{x},t)$ have the time-independent expectation value $\phi_n(\mathbf{x})$
\begin{equation}
\frac{\langle\Omega|\Phi_n(\mathbf{x},t)|\Omega\rangle}
{\langle\Omega|\Omega\rangle}
=\phi_n(\mathbf{x}).
\tag{16.3.4}\label{eq:16.3.4}
\end{equation}
It is convenient also to impose the condition that $\ket{\Omega}$ is normalized
\begin{equation}
\langle\Omega|\Omega\rangle=1.
\tag{16.3.5}\label{eq:16.3.5}
\end{equation}
To minimize the expectation value \eqref{eq:16.3.3} subject to the constraints \eqref{eq:16.3.4} and \eqref{eq:16.3.5}, we use the method of Lagrange multipliers, and instead minimize the quantity
\begin{equation}
\langle\Omega|H|\Omega\rangle
-\alpha\langle\Omega|\Omega\rangle
-\int d^3x\,\beta^n(\mathbf{x})
 \langle\Omega|\Phi_n(\mathbf{x})|\Omega\rangle
\tag{16.3.6}\label{eq:16.3.6}
\end{equation}
with no constraints on $\ket{\Omega}$. This gives
\begin{equation}
H\ket{\Omega}=\alpha\ket{\Omega}
+\int d^3x\,\beta^n(\mathbf{x})\Phi_n(\mathbf{x})\ket{\Omega}.
\tag{16.3.7}\label{eq:16.3.7}
\end{equation}
Both $\alpha$ and $\beta^n(\mathbf{x})$ are to be chosen in such a way as to satisfy the constraints \eqref{eq:16.3.4} and \eqref{eq:16.3.5} and therefore depend functionally on the prescribed expectation value $\phi_n(\mathbf{x})$.

Now, we have said that in the presence of a current $\mathcal J^n(\mathbf{x})$, the Hamiltonian $H-\int d^3x\,\mathcal J^n(\mathbf{x})\Phi_n(\mathbf{x})$ has an eigenvalue $E[\mathcal J]$:
\begin{equation}
\left[H-\int d^3x\,\mathcal J^n(\mathbf{x})\Phi_n(\mathbf{x})\right]
\ket{\Psi_{\mathcal J}}
=E[\mathcal J]\ket{\Psi_{\mathcal J}}
\tag{16.3.8}\label{eq:16.3.8}
\end{equation}
with a normalized eigenvector $\ket{\Psi_{\mathcal J}}$. Furthermore, since slowly turning on this current converts the vacuum into this energy eigenstate, we can presume that $E[\mathcal J]$ is the lowest energy eigenstate in the presence of this current. Therefore Eqs.~\eqref{eq:16.3.4}, \eqref{eq:16.3.5}, and \eqref{eq:16.3.7} are satisfied by
\begin{align}
\ket{\Omega}&=\ket{\Psi_{\mathcal J_\phi}},
&&\tag{16.3.9}\label{eq:16.3.9}\\
\alpha&=E[\mathcal J_\phi],
&&\tag{16.3.10}\label{eq:16.3.10}\\
\beta^n(\mathbf{x})&=\mathcal J_\phi^n(\mathbf{x}),
&&\tag{16.3.11}\label{eq:16.3.11}
\end{align}
where $\mathcal J_\phi(\mathbf{x})$ is the current for which $\Phi(\mathbf{x})$ has an expectation value $\phi(\mathbf{x})$ in the state $\ket{\Psi_{\mathcal J}}$.

Setting $\mathcal J=\mathcal J_\phi$ in Eq.~\eqref{eq:16.3.8} and taking the scalar product with $\ket{\Psi_{\mathcal J_\phi}}$, the minimum energy of states in which the fields $\Phi_n$ are constrained to have the expectation values $\phi_n$ is seen to be
\begin{equation}
\langle H\rangle_\Omega
=E[\mathcal J_\phi]
+\int d^3x\,\mathcal J_\phi^n(\mathbf{x})\phi_n(\mathbf{x}).
\tag{16.3.12}\label{eq:16.3.12}
\end{equation}
Recalling Eq.~\eqref{eq:16.3.2} and the assumed form of $J(x)$, this is
\begin{equation}
\langle H\rangle_\Omega
=\frac{1}{T}\left[-W[J_\phi]
+\int d^4x\,J_\phi^n(x)\phi_n(x)\right]
=-\frac{1}{T}\,\Gamma[\phi].
\tag{16.3.13}\label{eq:16.3.13}
\end{equation}
As noted in the previous section, if the field $\phi(x)$ has a constant value $\phi$ over a large spacetime volume $\mathcal V_4=\mathcal V_3T$, then we may write the effective action in terms of an effective potential $V(\phi)$:
\begin{equation}
\Gamma[\phi]=-\mathcal V_3T\,V(\phi).
\tag{16.3.14}\label{eq:16.3.14}
\end{equation}
In this case Eq.~\eqref{eq:16.3.13} tells us that the energy density is
\begin{equation}
\langle H\rangle_\Omega/\mathcal V_3=V(\phi).
\tag{16.3.15}\label{eq:16.3.15}
\end{equation}

This is the main result: $V(\phi)$ is \emph{the minimum of the expectation value of the energy density for all states constrained by the condition that the scalar fields $\Phi_n$ have expectation values $\phi_n$}. One consequence is that in the absence of external currents the vacuum state will relax to a state in which the potential $V(\phi)$ is not only stationary, which is required by the field equations \eqref{eq:16.1.8}, but also a minimum.

This result helps to resolve a problem in the interpretation of the quantum effective potential. The Euclidean version of the path-integral formalism (described in Appendix A of Chapter 23) makes it manifest that the two-point function $\Delta$ is positive (in the matrix sense), so according to Eq.~\eqref{eq:16.1.21} the same is true of $-\Pi=\Delta^{-1}$. Together with Eq.~\eqref{eq:16.2.3}, this implies that for a single scalar field the effective potential $V(\phi)$ must have a positive (or zero) second derivative with respect to $\phi$. More generally, the effective potential must be convex~\hyperref[ch16-ref-5]{[5]}:
\[
V\bigl(\lambda\phi_1+(1-\lambda)\phi_2\bigr)
\leq \lambda V(\phi_1)+(1-\lambda)V(\phi_2)
\quad\text{for }0\leq\lambda\leq1.
\]
But inspection of Eq.~\eqref{eq:16.2.13} shows that for $m^2<0$ and $g>0$ the zero-loop approximation to the effective potential for the scalar field theory with action \eqref{eq:16.2.1} has a \emph{negative-definite} second derivative when $\phi$ is between the two minima of the effective potential at $\pm\sqrt{6\lvert m^2\rvert/g}$. This contradiction arises because the derivation of perturbation theory implicitly relies on the existence of a stable vacuum, but when $V''(\phi)<0$, the field $\phi$ is at a value of $\bar\phi$ where $V(\bar\phi)-J_\phi\bar\phi$ is a maximum rather than a minimum, which means that the vacuum state in the presence of the current $J_\phi$ is unstable.

So what is the true effective potential for this scalar field theory when $m^2<0$ and $\phi$ lies between the two minima of the potential? The result of this section is that we must find the state of minimum energy in which the expectation value of the operator $\Phi$ equals $\phi$. As long as $\phi$ is between the two minima of the potential, we can give $\Phi$ an expectation value $\phi$ by taking the state as a suitable linear combination of the two states where $\phi$ is at the minima $\phi\simeq\pm\sqrt{6\lvert m^2\rvert/g}$. The energy in this state is equal to the energy at the minima, so this state clearly minimizes the energy. (Interference terms here vanish in the limit of infinite volume, for reasons explained in Section 19.1.) Thus the effective potential between the two minima of the potential is a \emph{constant}, satisfying the requirement that it have a nonnegative second derivative. The same argument shows that in more general theories where the potential has two local minima of unequal energy, the potential between these minima is linear~\hyperref[ch16-ref-6]{[6]}.
